\section{Literature Review}

\citet{fernandez2011} uses the co-movement of firms' financial
characteristics over time to define distances and converts those distances
into weights with an exponential kernel. By retaining her distance
while varying the weights, and then comparing
distances under a common weighting rule, we can trace differences in the
US return fits to the two choices her specification brings together.
Our empirical
criterion is the in-sample association of issuer returns with peer returns
after factor adjustment. The spatial asset-pricing theory of \citet{kou2018}
gives pricing implications under additional economic restrictions.

In financial representation learning, \citet{gu2021} use an autoencoder
asset-pricing model to learn how exposures vary with characteristics and to
estimate latent factors from returns. Tabular pretraining provides a
different starting point. Drawing on a model trained across datasets,
\citet{hollmann2025} demonstrate prediction and reusable embeddings with
TabPFN. We use a frozen encoder to obtain a vector for each issuer and
evaluate the financial meaning of the distances between those vectors
through the returns of the firms they bring together.

\citet{roeder2021} establish identification up to a linear transformation
for a family of discriminative models under their assumptions. A linear
transformation can stretch one direction more than another, changing the
ordering of distances even when the transformation is invertible. This
observation motivates our first proposition, which separates the information
a readout can recover from the choice of neighbours in a spatial model.
